\documentclass[11pt,a4paper]{article}
\usepackage[utf8]{inputenc}
\usepackage[T1]{fontenc}
\usepackage[margin=2cm]{geometry}
\usepackage{enumitem}
\usepackage{xcolor}
\usepackage{titlesec}
\usepackage{fancyhdr}
\usepackage{siunitx}

% Page style
\pagestyle{fancy}
\fancyhf{}
\fancyhead[L]{\small PhD Defense - Talking Points}
\fancyhead[R]{\small Page \thepage}
\renewcommand{\headrulewidth}{0.4pt}

% Section formatting
\titleformat{\section}{\Large\bfseries\color{blue!70!black}}{\thesection}{1em}{}
\titleformat{\subsection}{\large\bfseries\color{blue!50!black}}{\thesubsection}{1em}{}
\titleformat{\subsubsection}{\normalsize\bfseries}{\thesubsubsection}{1em}{}

% Compact lists
\setlist[itemize]{leftmargin=*,topsep=2pt,itemsep=1pt,parsep=0pt}

% Title
\title{\textbf{PhD Defense Presentation}\\[0.5em]
\Large Bolometry Diagnostics and Real-Time Radiation Feedback Control\\
at Wendelstein 7-X\\[1em]
\large Talking Points Guide}
\author{}
\date{10.04.2026}

\begin{document}

\section{TITLE SLIDE}

\begin{itemize}
\item Welcome to my PhD defense on bolometry diagnostics and real-time radiation feedback control at Wendelstein 7-X
\item This work addresses two critical challenges for future fusion reactors: (1) protecting plasma-facing components through controlled radiation, and (2) achieving stable detachment at high radiation fractions ($f_{\text{rad}}\geq85\%$)
\item Three main contributions: implemented first bolometer-based real-time feedback at W7-X, optimized channel selection through systematic sensitivity analysis, and validated MFR tomography with radially dependent anisotropy weighting
\end{itemize}

\section{CONTENTS}

\begin{itemize}
\item Presentation structure: (1) Motivation - fusion energy and W7-X challenges, (2) Bolometer diagnostic system - 94 channels with $\sim$5\,cm spatial resolution, (3) Real-time feedback - achieved $f_{\text{rad}}\geq85\%$ with factor $\geq2$ heat load reduction, (4) LOS sensitivity - identified optimal channel configurations, (5) Tomography - MFR with RDA for 2D radiation reconstruction, (6) Conclusions and future work
\end{itemize}

\section{MOTIVATION}

\subsection{Nuclear Fusion: Energy Challenge}

\begin{itemize}
\item Nuclear fusion through D-T reactions offers clean, abundant energy - deuterium from seawater, tritium bred from lithium
\item Lawson criterion: $n_e \cdot \tau_E \cdot T \geq 3\times10^{21}$\,keV$\cdot$s/m$^3$ defines conditions for net energy gain ($Q>1$)
\item Requires extreme conditions: $T\sim10$--20\,keV (100--200 million °C) and confinement time $\tau_E\sim1$--10\,s
\item Magnetic confinement uses strong fields (2.5\,T at W7-X) to contain hot plasma away from material walls
\item Stellarators like W7-X achieve steady-state operation through 3D-shaped coils, unlike pulsed tokamaks
\end{itemize}

\subsection{Wendelstein 7-X}

\begin{itemize}
\item W7-X is the world's largest and most advanced stellarator: $R=5.5$\,m, $a=0.53$\,m, $B=2.5$\,T, optimized for steady-state operation
\item Key challenge: managing heat loads on plasma-facing components during long pulses (up to 30 minutes planned)
\item Island divertor concept: 5 helical island chains intercept heat flux, connection lengths 100--1000\,m (10--100$\times$ longer than tokamaks)
\item Detachment strategy: radiate power in scrape-off layer and edge to reduce divertor heat loads from $\sim$10\,MW/m$^2$ to $<5$\,MW/m$^2$
\item Goal: achieve stable detachment at $f_{\text{rad}}\geq85$--90\% without plasma disruption while maintaining core performance
\item This requires real-time control of radiation distribution - motivation for bolometer feedback system
\end{itemize}

\subsection{Plasma Radiation}

\begin{itemize}
\item Plasma radiation mechanisms: (1) Bremsstrahlung (free-free transitions, $\sim T_e^{1/2}$), (2) Line radiation (bound-bound, temperature-dependent), (3) Recombination (free-bound)
\item Impurity seeding strategy: inject low-Z (N$_2$) or high-Z (Ne, Ar) gases to enhance edge radiation while avoiding core contamination
\item Carbon and oxygen are intrinsic impurities from wall interactions - carbon dominates by factor $10^2\times$ over oxygen at W7-X
\item Radiation cooling function $P_{\text{rad}}/(V\cdot n_e\cdot n_{\text{imp}})$ peaks at different temperatures for each impurity: C peaks $\sim$10--50\,eV (edge), O peaks $\sim$100\,eV
\item Transport in W7-X: neoclassical baseline with anomalous contributions from turbulence (ITG, TEM modes)
\item Anomalous transport increases diffusion by factor 10--100 over neoclassical, critical for impurity transport modeling
\item Radiation fraction $f_{\text{rad}} = P_{\text{rad}}/P_{\text{ECRH}}$ must be balanced: $f_{\text{rad}}<70\%$ insufficient for divertor protection, $f_{\text{rad}}>95\%$ risks core cooling and confinement degradation
\item Optimal operating window: $f_{\text{rad}}=85$--90\% achieves detachment while maintaining plasma stability
\end{itemize}

\section{CORE BOLOMETER AT W7-X}

\subsection{Bolometer Diagnostic (Hardware)}

\begin{itemize}
\item Bolometers measure total plasma radiation power (broadband, 0.2--600\,nm) through resistive heating of metal absorbers
\item Metal resistor design: thin gold film (1.5\,$\mu$m) on mica substrate, absorbs radiation $\rightarrow$ temperature increase $\rightarrow$ resistance change
\item Wheatstone bridge configuration: AC excitation (5\,V, kHz) eliminates 1/f noise, differential measurement doubles signal
\item Three camera arrays provide comprehensive plasma coverage: HBC (32 channels, horizontal), VBCl/r (2$\times$31 channels, vertical)
\item Total 94 channels installed, 61 functional in OP1.2b campaign (some damaged by ECRH stray radiation)
\item Each channel integrates radiation along line of sight - use geometry matrix $\mathbf{T}$ and etendue $K_M$ to reconstruct 2D distribution
\item Spatial resolution $\sim$5\,cm at magnetic axis, determined by etendue width and viewing geometry
\end{itemize}

\subsection{Performance}

\begin{itemize}
\item Calibration is critical for absolute power measurements - three methods used: (1) laser, (2) electrical ohmic heating, (3) in-situ during operation
\item In-situ ohmic heating calibration: apply $U_{\text{cal}}=1.2$--2.5\,V for 1.6\,s, measure current evolution $I(t)$, extract $R_M$, $\tau_M$, $\kappa_M$ from exponential decay
\item Calibration parameters (OP1.2b median): $R_M=0.987\pm0.011$\,$\Omega$, $\kappa_M=0.724\pm0.055$\,A$^2$, $\tau_M=110.57\pm4.93$\,ms, remarkably stable across 1182 experiments
\item Etendue $K_M$ quantifies geometric light collection efficiency: accounts for detector-aperture geometry, partial shadowing, angle-dependent transmission
\item Time resolution configurable: \{0.4, 0.8, 1.6, 3.2, 6.4, 12.8\}\,ms sampling available, 1.6\,ms used for feedback (balances speed vs noise)
\item Performance metrics: SNR$>1000$ at high radiation, noise $\sigma_{\Delta U}=0.339\pm12.586$\,$\mu$V, Gaussian error 1.6\,$\mu$W, drift $0.057\pm27.621$\,$\mu$V/s
\item Chord brightness $P_{\text{ch}} = P_M/K_M$ provides radial profile information - shows radiation distribution from core to SOL
\item Absorption efficiency near unity ($>95\%$) across 0.2--600\,nm wavelength range - truly broadband measurement
\end{itemize}

\section{REAL-TIME RADIATION FEEDBACK CONTROL}

\subsection{Feedback System (Island Divertor)}

\begin{itemize}
\item W7-X island divertor concept: 5 helical island chains ($m/n=5/5$) at plasma edge intercept heat and particle flux
\item Each island chain has upper and lower divertor modules - 10 modules total arranged helically around torus
\item Connection lengths $L_c=100$--1000\,m (10--100$\times$ longer than tokamaks) allow extensive radiation cooling along field lines
\item Standard magnetic configuration: $\iota/2\pi=1$ (rotational transform), creates 5-fold symmetric island structure
\item Gas injection system: 10 piezo valves for fast response ($<10$\,ms), thermal helium beam valves for feedback control
\item Valve locations strategically placed: HM1-5 modules (upper/lower pairs) allow targeted seeding into specific island chains
\item Island X-points are radiation hot spots - field lines converge, increasing local emissivity by factor 2--5$\times$
\item Bolometer cameras positioned to view these X-points - critical for feedback sensitivity
\item Field line tracing shows direct connection between gas injection points and divertor targets via island chains
\end{itemize}

\subsection{Feedback System (Implementation)}

\begin{itemize}
\item First bolometer-based real-time radiation feedback system at W7-X, operational during OP1.2b campaign (2018)
\item Hardware: NI 6321 DAQ card for fast acquisition, dedicated analog outputs for PID controller, parallel processing architecture
\item Two prediction proxies implemented: (1) $P_{\text{pred}}^{(1)}$ = geometry-weighted sum over selected channels $S$, (2) $P_{\text{pred}}^{(2)}$ = single channel voltage (dimensionless, minimal overhead)
\item Channel selection $S$ optimized through systematic analysis: VBC subset $S=\{61,65,73,78,86\}$ viewing separatrix and X-points achieves $\geq85\%$ accuracy
\item Alternative HBC selections also viable: $S=\{4,7,8,20,23,27\}$ provides similar performance with different viewing geometry
\item PID controller parameters: $K_p=0.5$--2.0 (proportional gain), $K_i=0.1$--0.5 (integral), $K_d=0$ (derivative not needed due to slow plasma response)
\item Setpoint: $f_{\text{rad}}$ target (typically 85--90\%), controller adjusts gas valve opening to maintain radiation fraction
\end{itemize}

\subsection{Feedback System (Latency)}

\begin{itemize}
\item System latency breakdown: acquisition 13.6\,ms (minimum, validated with laser diode), FIFO smoothing $\sim$8\,ms ($M=10$ samples $\times$ 1.6\,ms), processing $<1$\,ms
\item Total feedback loop latency: 150--400\,ms dominated by plasma response time (gas transport, ionization, radiation buildup)
\item FIFO moving average filter: reduces noise by factor $\sqrt{M}$ but adds latency $M\cdot\Delta t/2 \approx 8$\,ms for $M=10$
\item Trade-off: faster sampling ($\Delta t=0.8$\,ms) reduces latency but increases noise, 1.6\,ms optimal for feedback application
\item Benchmark test with laser diode: confirmed 13.6\,ms acquisition latency, validated timing chain from photon to digital signal
\item Future improvements: reduce acquisition to $<10$\,ms through hardware upgrades, implement predictive algorithms to compensate for plasma response delay
\item Comparison: Thomson scattering $\sim$50\,ms, interferometry $\sim$1\,ms, but bolometry measures the actual control variable (radiation)
\end{itemize}

\subsection{Experimental Achievements (XP20181010.32)}

\begin{itemize}
\item XP20181010.32 - breakthrough experiment achieving stable feedback-controlled detachment with hydrogen seeding
\item Discharge parameters: duration 9.2\,s, $P_{\text{ECRH}}=6.23$\,MW, standard magnetic configuration ($\iota/2\pi=1$), total input energy 55\,MJ
\item Key achievement: sustained $f_{\text{rad}}\geq85\%$ with peaks reaching 100\% ($P_{\text{rad}}=P_{\text{ECRH}}$), no plasma disruption or termination
\item Divertor heat load reduction: factor $\geq2\times$ from baseline $\sim$10\,MW/m$^2$ to $<5$\,MW/m$^2$, validated by infrared thermography
\item Detachment signatures: C$^{2+}$ filterscope emission appears at $f_{\text{rad}}\sim50\%$, intensifies and moves poloidally as $f_{\text{rad}}$ increases
\item Prediction accuracy: $P_{\text{pred}}^{(1)}$ tracks $P_{\text{rad}}$ within 10--15\% throughout discharge, validates 5-channel VBC subset $S=\{61,65,73,78,86\}$
\item Plasma stored energy $W_{\text{dia}}$ maintained at $\sim$0.8--1.0\,MJ despite high radiation - demonstrates compatibility with good confinement
\item This demonstrates reactor-relevant scenario: high $f_{\text{rad}}$ for component protection while maintaining plasma performance
\end{itemize}

\subsection{Experimental Achievements (Discussion)}

\begin{itemize}
\item Bolometer feedback advantages: (1) directly measures control variable (total radiation), (2) broadband - insensitive to impurity mix, (3) fast response (1.6\,ms sampling)
\item Comparison with alternatives: density feedback (DI) is indirect - affects confinement, can destabilize; C-III filterscope measures only one charge state, misses O, other C ions
\item Bolometer robustness: total radiation measurement independent of impurity species composition (C vs O vs seeded gases)
\item Spatial information: multiple viewing chords reveal radiation distribution - core vs edge vs SOL, identify X-point hot spots
\item Physical connection: radiation IS the detachment mechanism - controlling $P_{\text{rad}}$ directly controls heat dissipation
\item Gas injection optimization: moderate puff duration (50--200\,ms) optimal - short pulses cause transients, long pulses ineffective
\item PID tuning challenges: parameters discharge-specific, depend on magnetic configuration, heating power, impurity species
\item No universal scaling law found - complex parameter space requires adaptive control strategies
\item System limitations: 150--400\,ms total latency (dominated by plasma response), limits achievable control bandwidth to $\sim$2--5\,Hz
\item Future improvements: (1) reduce acquisition latency to $<10$\,ms, (2) implement predictive control using plasma models, (3) machine learning for automatic PID optimization
\item Despite limitations, bolometer feedback is the most direct method for radiation control in reactor-relevant scenarios
\item Validated approach: successfully demonstrated at W7-X, applicable to ITER, DEMO, and future stellarators
\end{itemize}

\section{LINE-OF-SIGHT SENSITIVITY}

\subsection{LOS Sensitivity (Methodology)}

\begin{itemize}
\item Systematic LOS sensitivity analysis to identify optimal channel configurations for feedback control
\item Dataset: 45 discharge pairs from OP1.2b campaign spanning wide parameter range ($P_{\text{ECRH}}=3.5$--6.5\,MW, $f_{\text{rad}}=0.3$--1.0, $T_e=0.9$--1.5\,keV)
\item Methodology: evaluate all possible channel combinations ($C(32,m)$ for HBC, $C(62,m)$ for VBC) for $m=3,5,7$ channels
\item Six evaluation metrics applied: (1) correlation, (2) weighted deviation, (3) mean deviation, (4) Fourier transform correlation, (5) Fisher information, (6) $\chi^2$ fitness
\item Each metric quantifies how well $P_{\text{pred}}$ from channel subset $S$ predicts total $P_{\text{rad}}$
\end{itemize}

\subsection{LOS Sensitivity (Results)}

\begin{itemize}
\item Key finding: channels viewing separatrix ($\rho_{\text{pol}}\sim0.95$--1.05) and SOL most effective - this is where detachment radiation concentrates
\item VBC cameras outperform HBC for edge sensitivity: better viewing geometry toward X-points and island chains
\item Global maximum in sensitivity profiles: VBC channel viewing upper inboard X-point ($r_{\text{eff}}\approx-0.13\,r_a$)
\item HBC shows local maxima at $\pm0.4\,r_a$, $\pm0.7\,r_a$ corresponding to island chain intersections
\item Robustness: multiple channel combinations achieve $\geq85\%$ prediction accuracy - no single ``best'' set exists
\item Optimal channel count: $m=3$--5 sufficient, $m=7$ provides marginal improvement ($<5\%$) at cost of increased complexity
\item Weighted deviation threshold: $\Delta_w<0.15$ defines acceptable performance (accounts for measurement uncertainties and temporal variations)
\item Correlation with plasma parameters: high $P_{\text{ECRH}}$ ($\geq4.2$\,MW), high $f_{\text{rad}}$ ($>0.75$), $T_e\sim1.3$\,keV yield best prediction quality
\end{itemize}

\section{TOMOGRAPHY}

\subsection{Minimum Fisher Regularization (Theory)}

\begin{itemize}
\item Tomography problem: invert line-integrated measurements ($\mathbf{T}\mathbf{x}=\mathbf{b}$) to reconstruct 2D emissivity distribution $\mathbf{x}$
\item Ill-posed inverse problem: 61 measurements but $\sim$4500 pixels (30 radial $\times$ 150 poloidal grid) - severely underdetermined
\item Geometry matrix $\mathbf{T}$: encodes etendue-weighted line integrals, condition number $\sim10^6$ indicates strong ill-posedness
\item Regularization required: Minimum Fisher Regularization (MFR) with Tikhonov-type penalty minimizes $\|\mathbf{T}\mathbf{x}-\mathbf{b}\|^2 + \lambda^{-1}\|\mathbf{H}\mathbf{x}\|^2$
\item Fisher information $I_F = \int(1/x)(\partial x/\partial r)^2\,dr$ measures gradient content - lower $I_F$ = smoother solution
\item Regularization matrix $\mathbf{H} \propto (1/x)\nabla^T\nabla$ penalizes large gradients weighted by local emissivity
\item Radially dependent anisotropy (RDA): different smoothness in radial vs poloidal directions, $k_{\text{ani}}(r)$ varies with radius
\item Core ($k_{\text{ani}}\approx2$--10): favors smooth, flux-surface aligned emission; Edge ($k_{\text{ani}}\approx0.1$--0.5): allows localized, anisotropic structures
\item Iterative solution: converges in $N_T=10$--20 iterations, regularization parameter $\lambda$ chosen by L-curve method
\item Computational grid: 30$\times$150 cells (radial$\times$poloidal), domain extends to $1.3^2V_P$ to handle boundary radiation
\end{itemize}

\subsection{Minimum Fisher Regularization (RDA)}

\begin{itemize}
\item RDA implementation: $k_{\text{ani}}(r)$ profile transitions smoothly from core to edge using tanh function
\item Typical profiles: $k_{\text{ani}}=\{2.0, 0.3\}$ (core, edge) for standard cases, $\{5.0, 0.5\}$ for highly anisotropic phantoms
\item Physical motivation: core radiation follows flux surfaces (isotropic in flux coordinates), edge shows X-point localization (anisotropic)
\item Advantages of MFR+RDA: (1) incorporates physics knowledge, (2) robust to noise, (3) stable convergence, (4) no negative emissivities
\item Alternative: Relative Gradient Smoothing (RGS) can be combined with RDA for enhanced feature localization
\item RGS weighting: $\mathbf{H}_{\text{RGS}} \propto (\nabla x/x)^2$ emphasizes relative gradients, stronger localization but may sacrifice global accuracy
\item Comparison: RDA chosen as primary method - better balance between accuracy and robustness for W7-X geometry
\end{itemize}

\subsection{Benchmarking (Geometry)}

\begin{itemize}
\item Geometry perturbation study: tested camera position variations ($\pm2^\circ$ tilts), detector/aperture segmentation ($N=2,4,8$), artificial symmetric geometries
\item Segmentation convergence: $N\geq8$ provides adequate resolution (etendue variations $<10^{-11}$\,mm$^{-3}$), $N=8$ chosen for computational efficiency
\item Camera tilt sensitivity: $\pm2^\circ$ rotations cause etendue changes $\sim10^{-11}$\,mm$^{-3}$, forward model variations 0.5--1\% in chord brightness
\item Intrinsic asymmetry: HBC 68.75$^\circ$ tilt causes 0.5--1\% left-right asymmetry in forward calculations even with symmetric radiation
\item Artificial symmetric HBC: mirroring half the array reduces asymmetry, validates that as-designed geometry is near-optimal
\item Worst-case scenario: $\pm2^\circ$ misalignment could shift reconstructed features by $\sim$0.2\,$r_a$, emphasizes need for accurate as-built measurements
\item Etendue validation: numerical integration agrees with analytical solutions for simple geometries (parallel plates, point sources)
\item Conclusion: discretization computationally stable, geometry uncertainties non-negligible but manageable with careful calibration
\end{itemize}

\subsection{Phantom Reconstructions}

\begin{itemize}
\item Phantom benchmark suite: 12+ synthetic radiation distributions testing algorithm performance across parameter space
\item Test cases include: (1) simple rings at various radii, (2) Gaussian features, (3) anisotropic island-like structures, (4) nested concentric profiles, (5) asymmetric distributions
\item Noise robustness: tested with up to 10\% Gaussian noise on measurements, algorithm maintains $\rho_c>0.85$ and $\chi^2<2$
\item Parameter variations: systematic scans of $k_{\text{ani}}$ (0.1--20), radial position (0.3--1.3\,$r_a$), intensity (0.1--2\,MW/m$^3$), width ($\sigma_r$, $\sigma_\theta$)
\item Quality metrics: (1) correlation $\rho_c$ (phantom vs tomogram), (2) $\chi^2$ fitness (forward vs measured), (3) MSD (mean squared deviation), (4) integrated powers $P_{\text{rad},2D}$
\item Key findings: $P_{\text{rad}}$ underestimates $P_{2D}$ by $\sim P_{2D}^{(\text{SOL})}$ (standard extrapolation misses SOL contribution), $\chi^2$ and $\rho_c$ not always congruent
\item Optimal parameters identified: $k_{\text{ani}}=\{2.0, 0.3\}$ for standard cases, $\{5.0, 0.5\}$ for highly anisotropic, $N_T=15$ iterations, grid 30$\times$150
\item Reconstruction quality: $>90\%$ for well-constrained phantoms (good camera coverage), 70--85\% for challenging cases (localized features, poor viewing angles)
\item Recommended settings validated through 59 phantom tests with comprehensive parameter variations
\end{itemize}

\subsection{Experimental Reconstructions}

\begin{itemize}
\item Experimental reconstructions: applied MFR+RDA to 4 feedback-controlled discharges from OP1.2b (XP20180809.13, XP20181010.32, XP20181011.12, XP20181018.45)
\item Temporal evolution: reconstructions show radiation moving from core to edge as $f_{\text{rad}}$ increases, validates detachment progression
\item Spatial features resolved: X-point localization (factor 2--3$\times$ local enhancement), island chain structure, MARFE-like asymmetries
\item XP20180809.13 example: laser blow-off events clearly visible as transient core radiation spikes, demonstrates temporal resolution
\item Power balance application: $P_{2D}^{(\text{core})} + P_{2D}^{(\text{SOL})}$ provides alternative to standard $P_{\text{rad}}$ extrapolation
\item Agreement: $P_{2D}^{(\text{total})}$ within 10--20\% of $P_{\text{rad}}$, discrepancy attributed to geometry uncertainties and SOL extrapolation errors
\item Key finding: standard extrapolation systematically underestimates total power by missing SOL contribution ($\sim$15--25\%)
\item Core radiation $P_{2D}^{(\text{core})}$ more stable metric for quasi-steady-state analysis - less sensitive to edge fluctuations
\item Statistical analysis (4 discharges, 50+ time slices): $P_{\text{rad}} \propto P_{2D}^{(X)}$ correlation strong ($R^2>0.85$), validates reconstruction accuracy
\item Quality metrics: $\chi^2=1.2$--2.5 (acceptable), $\rho_c=0.75$--0.90 (good correlation), MSD$<0.3$ (low deviation)
\item Discrepancy between $\chi^2$ and $\rho_c$: optimal tomogram may not have $\chi^2=1$ due to systematic errors (geometry, calibration)
\item Anisotropy optimization remains challenging: high-dimensional parameter space ($k_{\text{core}}$, $k_{\text{edge}}$, transition width), profile-dependent
\item Computational performance: $\sim$2--5 seconds per frame on standard CPU, suitable for post-shot analysis but not real-time
\item Future prospects: GPU acceleration could reduce to $<100$\,ms per frame, enabling real-time tomography for advanced feedback
\end{itemize}

\section{CONCLUSIONS}

\subsection{Key Achievements}

\begin{itemize}
\item Three major achievements form integrated diagnostic and control system for W7-X radiation management
\item (1) Real-time feedback system: first bolometer-based radiation control at W7-X, achieved $f_{\text{rad}}\geq85\%$ with factor $\geq2$ heat load reduction
\item (2) LOS optimization: systematic analysis of 45 discharge pairs, identified robust channel configurations ($m=3$--5 sufficient)
\item (3) Tomography: validated MFR+RDA algorithm through 59 phantom tests and 4 experimental reconstructions
\item Integration: feedback uses optimized channels, tomography provides detailed 2D validation of radiation distribution
\item Reactor relevance: demonstrated stable high-radiation scenarios essential for ITER, DEMO divertor protection
\end{itemize}

\textbf{Comprehensive Summary:}
\begin{itemize}
\item Feedback system specifications: minimum acquisition latency 13.6\,ms, total loop 150--400\,ms (plasma-limited), PID control with dual prediction proxies
\item Benchmark discharge XP20181010.32: 9.2\,s duration, $P_{\text{ECRH}}=6.23$\,MW, $f_{\text{rad}}\geq85\%$ sustained, heat load reduction factor $\geq2$, no disruption
\item Diagnostic characterization: 94 channels installed (61 functional OP1.2b), calibration stability $<5\%$ over 1182 experiments, SNR$>1000$
\item LOS sensitivity results: separatrix/SOL viewing channels optimal, VBC outperforms HBC for edge, multiple viable configurations identified
\item Detachment physics insights: C$^{2+}$ signature at $f_{\text{rad}}\sim50\%$, full detachment $f_{\text{rad}}>90\%$, radiation shifts inward to $\rho_{\text{pol}}\approx0.95$
\item STRAHL modeling: carbon dominates by $10^2\times$ over oxygen, C$^{3+}$/C$^{2+}$ critical near LCFS, validates feedback channel selection
\item Tomography performance: reconstruction quality $\rho_c>0.85$, $\chi^2<2.5$, $P_{2D}$ within 10--20\% of $P_{\text{rad}}$, resolves X-point structures
\item Power balance improvement: $P_{2D}^{(\text{core})}$ provides stable metric, reveals 15--25\% SOL underestimation in standard extrapolation
\item All methods extensively validated through OP1.2b campaign (2018): 1182 experiments, 45 feedback discharges analyzed
\end{itemize}

\section{OUTLOOK}

\subsection{Future and Possibilities}

\begin{itemize}
\item System upgrades roadmap: (1) hardware - reduce latency to $<10$\,ms through faster DAQ, (2) software - implement predictive control algorithms
\item Machine learning applications: (1) automatic PID optimization for different scenarios, (2) optimal channel selection, (3) predictive gas injection timing
\item Advanced modeling: integrate STRAHL with EMC3-EIRENE for 3D impurity transport, validate against OP2.1 high-performance scenarios
\item Improved chamber models: refine two/three-chamber parameters with OP2.1 data, develop predictive scaling laws for gas injection
\item Tomography extensions: (1) expand phantom benchmark suite to $>50$ cases, (2) GPU acceleration for real-time capability ($<100$\,ms/frame)
\item Real-time tomography vision: 2D radiation distribution feedback would enable spatial control (target specific regions), detect MARFEs early
\item Automated parameter optimization: Bayesian methods or grid search to find optimal $k_{\text{ani}}$ profiles for different radiation scenarios
\item ITER/DEMO relevance: methods directly transferable, DEMO requirements ($P_{\text{fus}}\sim500$\,MW, $f_{\text{rad}}\geq95\%$, $q_{\text{div}}<10$\,MW/m$^2$) demand reliable radiation control
\item Stellarator applicability: techniques applicable to other stellarators (LHD, future devices), tokamaks (ITER, DEMO, SPARC)
\item Long-term vision: fully automated radiation control system with adaptive learning, spatial distribution optimization, predictive capabilities
\end{itemize}

\section{CLOSING}

\subsection{Thank You}

\begin{itemize}
\item Thank you for your attention - ``Sin é, a chairde'' (That's it, friends) - Irish Gaelic closing
\item Open for questions on any aspect: feedback system implementation, LOS optimization methodology, tomography algorithm details, experimental results interpretation
\item Key discussion points prepared: (1) latency reduction strategies, (2) comparison with alternative feedback methods, (3) tomography parameter optimization, (4) ITER/DEMO applicability
\item Acknowledge W7-X team, IPP collaborators, EUROfusion support, Helmholtz Association, University of Greifswald
\end{itemize}

\newpage
\section*{BACKUP SLIDES - Quick Reference}

\textbf{Plasma Physics:}
\begin{itemize}
\item Transport mechanisms: classical, neoclassical, anomalous
\item Impurity transport: D$_{\text{an}}\sim0.3$--3\,m$^2$/s dominates
\item Detachment: T$_e < 5$\,eV threshold, P$_{\text{div}} \rightarrow 0$
\end{itemize}

\textbf{Modeling \& Analysis:}
\begin{itemize}
\item Two/three-chamber seeding models
\item Power balance: P$_{\text{bal}} \approx 0$ within errors
\end{itemize}

\textbf{Diagnostic Details:}
\begin{itemize}
\item Etendue: $\widetilde{K}_M$ quantifies light collection
\item Calibration: R$_M$, $\tau_M$, $\kappa_M$ from ohmic heating
\item Performance: SNR$>1000$, noise 0.339\,$\mu$V
\item Bolometer equation: AC bridge, power balance
\end{itemize}

\textbf{STRAHL Modeling:}
\begin{itemize}
\item 1D impurity transport code
\item Carbon dominance: $10^2\times$ over oxygen
\item Critical charge states: C$^{3+}$/C$^{2+}$ near LCFS
\item Radiation shift with f$_{\text{rad}}$
\end{itemize}

\textbf{LOS Sensitivity Details:}
\begin{itemize}
\item Six metrics: correlation, deviation, FFT, Fisher, $\chi^2$
\item Geometry sensitivity: $\pm2^\circ$ tilt acceptable
\item Segmentation: N=8 optimal
\item RDA vs RGS: RDA more robust
\item Plasma parameters: P$_{\text{ECRH}}\geq4.2$\,MW optimal
\end{itemize}

\textbf{DAQ Latency:}
\begin{itemize}
\item Laser characterization: 13.6\,ms minimum
\item Sample times: 0.8--3.2\,ms tested
\item Total loop: 150--400\,ms sufficient
\end{itemize}

\vfill
\noindent\textit{Document prepared for PhD defense presentation}\\
\textit{Use this guide to maintain consistent messaging and technical accuracy}

\end{document}

% Made with Bob
